%%% FIELD cidp-contract-narrowed-statement
When \(S=\varnothing\), Algorithm~2 CIDP exchanges only nonempty intersections of the current intervention or conditioning set with a bucket, and performs each exchange only after the corresponding displayed PAG Rule--2 separation test succeeds.  On success it reaches a role state \(\rho_\star=(\bar X_\star,\bar C_\star)\), its final marginal-algorithm call returns an observed-law expression for \(P_{\mathbf M}(Y,\bar C_\star\mid\operatorname{do}(\bar X_\star))\), and CIDP divides that expression by its sum over the full discrete support of \(Y\).  Theorem~4 states the completeness direction for identification of \(P_{\mathbf M}(Y\mid\operatorname{do}(X),C)\) from an input-free PAG: if this conditional effect is identifiable, CIDP does not return \(\operatorname{FAIL}\).
%%% FIELD no-selection-decision-statement
When \(S=\varnothing\), \(\operatorname{Branch\text{-}sCIDP}(\mathcal P;Y,X,C)\) and \(\operatorname{CIDP}(\mathcal P;Y,X,C)\) agree on the identification decision for
\[
 P_{\mathbf M}(Y=y\mid\operatorname{do}(X=x),C=c);
\]
that is, one returns a formula if and only if the other returns a formula.
%%% FIELD no-selection-decision-proof
Assume \(S=\varnothing\).  Since the selection event is then certain, Definition \(\mathrm{def{:}target\text{-}query}\) gives
\[
 \psi_{\mathbf M}(y;x,c)
 =P_{\mathbf M}(Y=y\mid\operatorname{do}(X=x),C=c).
 \tag{1}
\]
The conditioning stratum in (1) is positive by \(\mathrm{ass{:}positive\text{-}discrete\text{-}selected\text{-}model}\), as incorporated into that definition.

First suppose that Branch-sCIDP returns a formula.  By Theorem \(\mathrm{thm{:}branch\text{-}soundness}\), that formula equals (1) for every model \(\mathbf M\in\mathfrak M_d^+(\mathcal P)\).  Thus the conditional effect in (1) is identified from the represented observational law.  The completeness direction in Lemma \(\mathrm{lem{:}cidp\text{-}no\text{-}selection\text{-}contract}\) says that CIDP cannot return \(\operatorname{FAIL}\) for an identifiable conditional effect.  Hence CIDP returns a formula.

Conversely, suppose that CIDP returns a formula.  By the transition contract in Lemma \(\mathrm{lem{:}cidp\text{-}no\text{-}selection\text{-}contract}\), its successful execution produces a sequence
\[
 (X,C)=\rho_0\longrightarrow\rho_1\longrightarrow\cdots
 \longrightarrow\rho_k=(\bar X_\star,\bar C_\star),
 \tag{2}
\]
where every arrow moves a nonempty intersection of a current role set with a bucket and is taken only after the corresponding PAG Rule--2 separation test succeeds.  These are exactly the edge predicates in Definition \(\mathrm{def{:}reachable\text{-}role\text{-}graph}\).  Consequently
\[
 \rho_\star=(\bar X_\star,\bar C_\star)
 \in\Omega(\mathcal P;Y,X,C).
 \tag{3}
\]
The same cited contract says that the final marginal call succeeds with an observed-law expression \(E\) for
\[
 Q_{\rho_\star}
 :=P_{\mathbf M}(Y,\bar C_\star
       \mid\operatorname{do}(\bar X_\star)).
 \tag{4}
\]
Thus (4) is identified on the represented PAG model class and, in particular, on its positive discrete subclass \(\mathfrak M_d^+(\mathcal P)\).

The hypotheses of Lemma \(\mathrm{lem{:}sidp\text{-}completeness}\) apply to the terminal joint query with
\[
 A_0=Y\cup\bar C_\star,
 \qquad B_0=\bar X_\star.
\]
Indeed, \(A_0\neq\varnothing\) by \(\mathrm{ass{:}nonempty\text{-}outcome}\).  Moreover, \(A_0\cap B_0=\varnothing\): Definition \(\mathrm{def{:}reachable\text{-}role\text{-}graph}\) keeps \(\bar X_\star\mathbin{\dot\cup}\bar C_\star=X\cup C\), while \(\mathrm{ass{:}disjoint\text{-}query\text{-}sets}\) gives \(Y\cap(X\cup C)=\varnothing\).  If the sIDP call at \(\rho_\star\) returned \(\operatorname{Fail}\), Lemma \(\mathrm{lem{:}sidp\text{-}completeness}\) would say that (4) is not identified over \(\mathfrak M_d^+(\mathcal P)\), contradicting the preceding paragraph.  Hence that sIDP call succeeds.

By Definition \(\mathrm{def{:}branch\text{-}scidp}\), Branch-sCIDP enumerates every state in \(\Omega(\mathcal P;Y,X,C)\) and calls sIDP at each state.  In view of (3), it therefore succeeds either before reaching \(\rho_\star\) or at \(\rho_\star\).  The normalization at the latter state is defined: positivity gives
\[
 \sum_{y'}E(y',\bar c_\star;\bar x_\star)
 =P_{\mathbf M}(\bar C_\star=\bar c_\star
       \mid\operatorname{do}(\bar X_\star=\bar x_\star))>0.
\]
The sum is over the discrete outcome support; if that support is countably infinite, the equality is countable additivity of nonnegative point masses.  When \(\bar C_\star=\varnothing\), the same sum equals one.  We have proved both implications between the two success decisions, which is the claimed equivalence.
